\chapter{Conclusion}\label{conclusion}

[Broadly, my goal to create a functional web app was not met]

[But to be fair, it was ambitious -- Last.fm web apps are usually much simpler than this, and using scraping in this way would be unique, which is obviously why it's presenting the most challenge in find a suitable deployment option]

[There's no time crunch to get a site like this deployed, and the community is such that I can ask for advice on topics like web app deployment to see how applicable others' experiences would be to my situation]

Broadly, the overall goal of this project has been to create a web application using web scraping to access Last.fm data on the size of artists' listenerships in order to inform a user how mainstream their recent music listening has been. In the strictest sense, this goal has not been met -- the application is not deployed, and users cannot access it. Still, this certainly does not mean the effort put into it was for naught, or that the endeavor was a total failure. Indeed, the software functions properly na efficiently as a proof-of-concept and much of the functionality present in the original version of the software -- a simple text-based command line tool -- was replicated or expanded on in the web version, giving further indication of the project's potential.

The extent of tools and frameworks available for use in web applications allowed the project to grow in scope. The ability of frontend tools and frameworks like JavaScript and D3 allowed for the visualization of data in such a way that a user gets a much better idea of the shape of their data than could've been provided by the purely text-based conclusions the data analysis could determine. In addition, the ease of connecting the backend to a data store allowed for a method of caching scraped data to save time on redundant future processes. These processes, common and natural among web applications, were able to bring depth and utility to the existing project.

Even more web-based utilities will be of use in future development of the project. In addition to the APIs connecting the components of the software, gaining access to the official API will allow it to communicate with Last.fm in an official and intentional capacity while also reducing its dependencies by no longer relying on other external tools.

The limitations in the finished project are not strictly idea, but they are not wholly disappointing, either. Compared to other Last.fm-based web apps, this one is notably ambitious. Going without the API and instead relying on scraping -- a fairly intensive and time-consuming process to delegate to a backend is not implemented in any other tools, as far as I know. A conscious choice was made to de-prioritize the research required for finding a suitable hosting service, instead opting to refine the software and improve it as a proof-of-concept. Thus far, that choice feels like a sensible one.

There is no time limit on a project like this. Last.fm and the web apps that rely on and augment it have been around for decades now, and will likely remain for years into the future. In that time, more progress can be made on the software and in researching and improving at the skills and knowledge used to create it. Other open-source developers are out there making and maintaining different tools, and their experience can be learned from both impersonally, by studying their software; and personally by reaching out and asking questions. This project serves as the start of a journey to step into this ecosystem, which I hope will continue for me long after it concludes.